\cleardoublepage
\vspace{2cm}
\begin{center}
{\fontsize{14}{16.8}\selectfont\textbf{DAFTAR SIMBOL}}\\[1em]
\end{center}
\label{chap:simbol}
\addcontentsline{toc}{chapter}{DAFTAR SIMBOL}
\setcounter{chapter}{0}
\setcounter{section}{0}
\vspace{2em}

\begin{longtable}{@{} l >{\raggedright\arraybackslash}p{7cm} c @{}}
\multicolumn{1}{l}{\textbf{Notasi}} & \multicolumn{1}{l}{\textbf{Deskripsi}} & \multicolumn{1}{l}{\textbf{Hal.}} \\[0.5ex]

$s$                              & Indeks mahasiswa                                                              & \pageref{chap:perancangan} \\
$p$                              & Indeks kegiatan kolaboratif (\emph{project})                                  & \pageref{chap:perancangan} \\
$i$                              & Indeks peringkat atribut dalam sub-algoritma                                  & \pageref{chap:studi-literatur} \\
$k$                              & Indeks penjumlahan pada formula \emph{Rank Order Centroid}                    & \pageref{chap:studi-literatur} \\
$n$                              & Jumlah atribut dalam satu sub-algoritma                                       & \pageref{chap:studi-literatur} \\
$w_i$                            & Bobot \emph{Rank Order Centroid} untuk atribut pada peringkat ke-$i$          & \pageref{chap:studi-literatur} \\
$K_s$                            & Himpunan \emph{skill} yang dimiliki mahasiswa $s$                             & \pageref{chap:perancangan} \\
$K_p$                            & Himpunan \emph{skill} yang dibutuhkan kegiatan $p$                            & \pageref{chap:perancangan} \\
$M_s$                            & Himpunan tag minat yang dimiliki mahasiswa $s$                                & \pageref{chap:perancangan} \\
$M_p$                            & Himpunan tag minat yang relevan dengan kegiatan $p$                           & \pageref{chap:perancangan} \\
$\text{expLevel}_s$              & Level pengalaman mahasiswa $s$ (skala ordinal 1--3)                           & \pageref{chap:perancangan} \\
$\text{expReq}_p$                & Level pengalaman yang dibutuhkan kegiatan $p$ (skala ordinal 1--3)            & \pageref{chap:perancangan} \\
$\text{gayaKerja}_s$             & Gaya kerja mahasiswa $s$ (Terstruktur atau Fleksibel)                         & \pageref{chap:perancangan} \\
$\text{gayaKerjaPreferensi}_p$   & Preferensi gaya kerja yang diinginkan kegiatan $p$                            & \pageref{chap:perancangan} \\
$\text{preferensiPeran}_s$       & Preferensi peran mahasiswa $s$                                                & \pageref{chap:perancangan} \\
$\text{rolesDibutuhkan}_p$       & Himpunan peran yang dibutuhkan kegiatan $p$                                   & \pageref{chap:perancangan} \\
$\text{SkorSkill}(s,p)$          & Skor kecocokan \emph{skill} antara mahasiswa $s$ dan kegiatan $p$             & \pageref{chap:perancangan} \\
$\text{SkorPengalaman}(s,p)$     & Skor kecocokan pengalaman antara mahasiswa $s$ dan kegiatan $p$               & \pageref{chap:perancangan} \\
$\text{SkorMinat}(s,p)$          & Skor kecocokan minat antara mahasiswa $s$ dan kegiatan $p$                    & \pageref{chap:perancangan} \\
$\text{SkorGayaKerja}(s,p)$      & Skor kecocokan gaya kerja antara mahasiswa $s$ dan kegiatan $p$               & \pageref{chap:perancangan} \\
$\text{SkorPeran}(s,p)$          & Skor kecocokan preferensi peran antara mahasiswa $s$ dan kegiatan $p$         & \pageref{chap:perancangan} \\
$\text{SkorAlgoritma1}(s,p)$     & Skor sub-algoritma 1 (\emph{Demands-Abilities Fit})                           & \pageref{chap:perancangan} \\
$\text{SkorAlgoritma2}(s,p)$     & Skor sub-algoritma 2 (\emph{Needs-Supplies Fit})                              & \pageref{chap:perancangan} \\
$\text{AffinityScore}(s,p)$      & Skor kecocokan akhir antara mahasiswa $s$ dan kegiatan $p$                    & \pageref{chap:perancangan} \\
$|\cdot|$                        & Kardinalitas himpunan atau nilai mutlak                                       & \pageref{chap:perancangan} \\
$\cap$                           & Irisan dua himpunan                                                           & \pageref{chap:perancangan} \\
$\in$                            & Keanggotaan dalam suatu himpunan                                              & \pageref{chap:perancangan} \\
$\sum$                           & Operator penjumlahan                                                          & \pageref{chap:studi-literatur} \\

\end{longtable}
